\chapter{Exercícios práticos}
\label{cap:exercicios}

Os exercícios seguem a ordem da apostila. As partes 1, 4 e 5 e o exercício 6.1 funcionam só com o Git instalado, sem internet e sem servidor. A parte 2 precisa de acesso a um GitLab. A parte 3 precisa de um agente de IA e de um projeto que você possa usar. Os exercícios 6.2 e 6.3 precisam de uma conta no GitHub ou no GitLab.

\begin{dica}
Faça os exercícios numa pasta de treino, fora de qualquer projeto real. Se algo der errado, apague a pasta e comece de novo: é para isso que ela existe.
\end{dica}

% -----------------------------------------------------------------------------
\section{Parte 1 — Git no seu computador}

\begin{exercicio}{1.1 — Primeiro repositório e primeiro commit}
\begin{enumerate}
  \item Crie um repositório local.
  \item Crie um arquivo.
  \item Execute \cmd{git add}.
  \item Execute \cmd{git commit}.
\end{enumerate}
\end{exercicio}

\begin{terminal}
mkdir exercicio-git
cd exercicio-git
git init -b main
echo "Olá, Git!" > saudacao.txt
git status
git add saudacao.txt
git commit -m "Cria a saudação"
\end{terminal}

\begin{verifique}
\cmd{git log --oneline} mostra uma linha com \enquote{Cria a saudação}, e \cmd{git status} diz que não há nada a submeter.
\end{verifique}

\begin{exercicio}{1.2 — Branch, alteração e merge}
\begin{enumerate}\setcounter{enumi}{4}
  \item Crie uma branch.
  \item Altere o arquivo.
  \item Crie outro commit.
  \item Faça o merge na \arq{main}.
\end{enumerate}
\end{exercicio}

\begin{terminal}
git switch -c feature/despedida
echo "Até logo!" >> saudacao.txt       # >> acrescenta uma linha no fim
git commit -am "Adiciona a despedida"  # -a prepara os arquivos já acompanhados
git switch main
git merge feature/despedida
cat saudacao.txt
\end{terminal}

\begin{verifique}
O merge diz \enquote{Fast-forward} e o arquivo tem duas linhas: \enquote{Olá, Git!} e \enquote{Até logo!}.
\end{verifique}

\begin{exercicio}{1.3 — Simular e resolver um conflito}
\begin{enumerate}\setcounter{enumi}{8}
  \item Simule um conflito.
  \item Resolva o conflito.
\end{enumerate}
A ideia é mudar \textbf{a mesma linha} de formas diferentes em duas branches.
\end{exercicio}

\begin{terminal}
# na branch nova, a saudação fica formal
git switch -c feature/formal
printf "Prezados, bom dia!\nAté logo!\n" > saudacao.txt
git commit -am "Deixa a saudação formal"

# na main, a mesma linha fica informal
git switch main
printf "E aí, pessoal!\nAté logo!\n" > saudacao.txt
git commit -am "Deixa a saudação informal"

# agora junte as duas
git merge feature/formal
\end{terminal}
\begin{saida}
Mesclagem automática de saudacao.txt
CONFLITO (conteúdo): conflito de mesclagem em saudacao.txt
Automatic merge failed; fix conflicts and then commit the result.
\end{saida}

Abra \arq{saudacao.txt} num editor de texto. Você vai ver os marcadores explicados na seção \ref{sec:conflitos}. Escolha a versão que fica (ou escreva uma terceira), apague os marcadores, salve e registre:

\begin{terminal}
git add saudacao.txt
git commit -m "Resolve o conflito da saudação"
git log --oneline --graph --all
\end{terminal}
\begin{saida}
*   757aae3 Resolve o conflito da saudação
|\
| * c2524a8 Deixa a saudação formal
* | a41aded Deixa a saudação informal
|/
* afe16bb Adiciona a despedida
* 4203776 Cria a saudação
\end{saida}

\begin{verifique}
O desenho do \cmd{git log --graph} mostra as duas linhas se separando e se juntando de novo no commit de resolução. Os códigos dos commits no seu computador serão outros.
\end{verifique}

% -----------------------------------------------------------------------------
\section{Parte 2 — GitLab}

\begin{exercicio}{2.1 — Um Merge Request completo}
\begin{enumerate}
  \item Crie ou use um projeto no GitLab.
  \item Crie uma branch.
  \item Faça uma alteração.
  \item Faça \cmd{push}.
  \item Crie um Merge Request.
  \item Revise a alteração.
  \item Faça o merge.
\end{enumerate}
\end{exercicio}

\begin{terminal}
git clone <URL_DO_PROJETO_DE_TREINO>
cd <pasta-do-projeto>
git switch -c exercicio/meu-nome
echo "Treinamento de Git na Marinha" >> README.md
git add README.md
git commit -m "Registra o treinamento no README"
git push origin exercicio/meu-nome
\end{terminal}

Depois, no navegador:

\begin{enumerate}
  \item abra o projeto e clique em \textbf{Create merge request} (ou vá em \textbf{Merge requests \faChevronRight\ New merge request});
  \item confira: origem \arq{exercicio/meu-nome}, destino \arq{main};
  \item escreva um título e uma descrição e crie o MR;
  \item abra a aba \textbf{Changes} e deixe um comentário numa linha;
  \item peça a um colega para aprovar e faça o merge (se o projeto permitir).
\end{enumerate}

\begin{dica}[Pelo terminal]
Os passos no navegador podem ser trocados por \cmd{glab mr create -b main}, \cmd{glab mr view} e \cmd{glab mr merge} (capítulo \ref{cap:glab}).
\end{dica}

\begin{verifique}
O MR aparece como \enquote{Merged}, e a linha nova está no \arq{README.md} da \arq{main}.
\end{verifique}

% -----------------------------------------------------------------------------
\section{Parte 3 — IA}

\begin{atencao}[Escolha bem o projeto]
Use um projeto pequeno, \textbf{pessoal ou de código aberto}. Não use código institucional sem autorização.
\end{atencao}

\begin{exercicio}{3.1 — Descoberta com um agente de IA}
Escolha um pequeno projeto existente e peça ao agente:
\begin{enumerate}
  \item realizar descoberta somente leitura;
  \item identificar tecnologias;
  \item identificar a arquitetura;
  \item identificar dependências;
  \item identificar funcionalidades;
  \item identificar problemas;
  \item gerar documentação;
  \item apresentar dúvidas e pontos que exigem validação humana.
\end{enumerate}
\end{exercicio}

Um ponto de partida é o prompt abaixo. Ele é uma \textbf{versão reduzida}, escrita para este exercício, que aproveita as travas do prompt final: somente leitura, nunca inventar, parar e perguntar.

\begin{terminalt}{· PROMPT DE EXEMPLO}
Você vai analisar este projeto. Siga estas regras sem exceção:

1. MODO SOMENTE LEITURA: não altere, não apague e não renomeie nenhum
   arquivo, e não faça commits. Você só pode criar arquivos em docs/.
2. NUNCA INVENTE: todo fato precisa de evidência (arquivo e linha).
   O que não puder confirmar, marque como DESCONHECIDO, INFERIDO ou
   REQUER VALIDAÇÃO.
3. PARE E PERGUNTE: se encontrar contradição, senha, token, dado pessoal
   ou algo que exija decisão humana, pare e me relate antes de seguir.
4. IDIOMA: escreva tudo em português brasileiro.

Tarefa:
a) Crie docs/DISCOVERY.md com a estrutura de pastas, as tecnologias e
   as dependências encontradas.
b) Crie docs/ARQUITETURA.md explicando os componentes e como se
   comunicam, com um diagrama em texto.
c) Crie docs/FUNCIONALIDADES.md listando o que o sistema faz.
d) Crie docs/PROBLEMAS.md com defeitos, riscos e dívida técnica.
e) Termine com uma lista de perguntas que eu preciso validar.
\end{terminalt}

\begin{exercicio}{3.2 — Validar o que a IA produziu}
Pegue três afirmações da documentação gerada e confira cada uma no código. Anote: estava certa, errada ou incompleta? Responda também às perguntas que o agente deixou no item (e).
\end{exercicio}

\begin{verifique}
\cmd{git status} mostra \textbf{só} arquivos novos em \arq{docs/}. Se aparecer qualquer outro arquivo modificado, o agente desrespeitou a regra 1: use \cmd{git diff} para ver o que mudou e \cmd{git restore} para desfazer.
\end{verifique}

% -----------------------------------------------------------------------------
\section{Parte 4 — Desafio: reproduzir o Problema 1}

\begin{exercicio}{4.1 — O caso real, sem servidor}
Simule o GitLab institucional com um repositório \enquote{vazio de trabalho} (\emph{bare}) numa pasta do seu computador, com a \arq{main} e o README padrão. Depois, integre a ele um sistema que já tem histórico próprio, seguindo os oito passos do capítulo \ref{cap:problema1}.
\end{exercicio}

\textbf{Preparação: o \enquote{GitLab} de mentira.}

\begin{terminal}
mkdir desafio && cd desafio
git init --bare -b main gitlab.git       # faz o papel do servidor
git clone gitlab.git instituicao
cd instituicao
echo "# Projeto (README padrão da instituição)" > README.md
git add README.md
git commit -m "Initial commit"
git push origin main
cd ..
\end{terminal}

\textbf{Preparação: o sistema que já existia.}

\begin{terminal}
mkdir sistema && cd sistema
git init -b main
echo "# Sistema de Controle" > README.md
echo "print('ola')" > app.py
git add README.md app.py
git commit -m "Primeira versão"
echo "x = 1" >> app.py
git commit -am "Ajuste"
\end{terminal}

\textbf{Agora é com você.} Ainda dentro de \arq{sistema/}, aplique os oito passos. Como o \enquote{servidor} é uma pasta, a URL é o caminho dela: \arq{../gitlab.git}. O passo 8 (Merge Request) não existe numa pasta, então pare no passo 7.

\begin{enumerate}
  \item Antes do passo 4, tente \cmd{git merge origin/main} sem o parâmetro. Qual mensagem aparece? Por quê?
  \item No passo 5, mantenha o README do \textbf{sistema}. Qual comando você usou: \cmd{--ours} ou \cmd{--theirs}?
  \item Depois do passo 7, rode \cmd{git log --oneline --graph --all}. Onde está o commit que une os dois históricos?
\end{enumerate}

\begin{tcolorbox}[base, colback=cinzaSuave, colframe=cinzaSuave, title={\faKey\enspace Respostas}, coltitle=azulMB, attach title to upper={\\[2pt]}]
\begin{enumerate}
  \item \texttt{fatal: refusing to merge unrelated histories}. Os dois repositórios nasceram separados e não têm ancestral comum.
  \item \cmd{git restore --ours README.md}: durante o merge você está na branch \arq{desenvolvimento}, que tem o README do sistema.
  \item No topo: o commit \enquote{Merge remote-tracking branch 'origin/main' into desenvolvimento}, com duas linhas chegando nele. Uma traz \enquote{Primeira versão} e \enquote{Ajuste}, a outra traz \enquote{Initial commit}.
\end{enumerate}
\end{tcolorbox}

% -----------------------------------------------------------------------------
\section{Parte 5 — Padrões de commit}

\begin{exercicio}{5.1 — Reescrever mensagens ruins}
Reescreva cada mensagem abaixo no formato Conventional Commits (capítulo \ref{cap:padroes}). Escolha o tipo, invente um escopo que faça sentido e escreva uma descrição curta.
\begin{enumerate}
  \item \enquote{arrumei o erro que dava quando a senha tava vazia}
  \item \enquote{coloquei explicação de como instalar no readme}
  \item \enquote{mudei a identação do arquivo de rotas}
  \item \enquote{agora o relatório sai em pdf também}
  \item \enquote{atualizei o pipeline pra rodar os testes}
  \item \enquote{a API agora exige token em tudo, quem não mandar leva 401}
\end{enumerate}
\end{exercicio}

\begin{exercicio}{5.2 — Aplicar no seu repositório de treino}
No repositório \arq{exercicio-git} da Parte 1, faça três commits pequenos, cada um com um tipo diferente (\arq{feat}, \arq{docs} e \arq{fix}). Depois rode \cmd{git log --oneline} e confira se dá para entender o histórico só lendo as mensagens.
\end{exercicio}

\begin{tcolorbox}[base, colback=cinzaSuave, colframe=cinzaSuave, title={\faKey\enspace Respostas possíveis do 5.1}, coltitle=azulMB, attach title to upper={\\[2pt]}]
\ttfamily\small
\begin{enumerate}
  \item fix(login): aceita tentativa com senha vazia sem erro
  \item docs(readme): explica a instalação
  \item style(rotas): corrige a indentação
  \item feat(relatorio): exporta em PDF
  \item ci: roda os testes no pipeline
  \item feat(api)!: exige token em todos os endpoints
\end{enumerate}
\rmfamily\normalsize
No item 6, o \arq{!} marca a quebra de compatibilidade. Um rodapé \arq{BREAKING CHANGE: clientes sem token recebem 401} deixaria isso ainda mais claro. Os escopos e as palavras podem variar: o importante é o tipo certo e uma descrição que diga o que mudou.
\end{tcolorbox}

% -----------------------------------------------------------------------------
\section{Parte 6 — Assinatura de commits}

\begin{exercicio}{6.1 — Ver o problema com os próprios olhos}
Num repositório de treino, faça um commit com um nome e um e-mail que não são seus e confira o que o \cmd{git log} mostra:
\end{exercicio}

\begin{terminal}
git -c user.name="Pessoa Qualquer" -c user.email="qualquer@exemplo.com" \
    commit --allow-empty -m "Quem fez este commit?"
git log --format='%an <%ae> | assinatura: %G?' -1
\end{terminal}
\begin{saida}
Pessoa Qualquer <qualquer@exemplo.com> | assinatura: N
\end{saida}

O Git aceitou a identidade falsa, e o \arq{N} diz que não há assinatura.

\begin{exercicio}{6.2 — Assinar de verdade}
Siga o capítulo \ref{cap:assinatura} com o método que preferir (GPG ou SSH):
\begin{enumerate}
  \item crie ou reaproveite uma chave;
  \item cadastre a chave no GitHub ou no GitLab;
  \item configure o Git para assinar todo commit;
  \item faça um commit num repositório seu e envie com \cmd{git push};
  \item confira o selo na página de commits.
\end{enumerate}
\end{exercicio}

\begin{verifique}
No terminal, \cmd{git log --format='\%h \%G? \%s' -1} mostra a letra \arq{G}. No site, o commit tem o selo \textbf{Verified}. Se aparecer \textbf{Unverified}, confira primeiro o e-mail: \cmd{git config user.email} precisa ser um e-mail verificado na sua conta.
\end{verifique}

\begin{exercicio}{6.3 — Uma tag assinada}
Crie e confira uma tag assinada:
\end{exercicio}

\begin{terminal}
git tag -s v1.0 -m "Versão 1.0"
git tag -v v1.0
git push origin v1.0
\end{terminal}

\begin{verifique}
O \cmd{git tag -v} mostra a mesma confirmação de assinatura do \cmd{git log --show-signature}, e a tag aparece como \textbf{Verified} na página de tags do GitHub.
\end{verifique}
